\section{Essentially constant systems}
\label{section-essentially-constant}

\noindent
Let $M : \mathcal{I} \to \mathcal{C}$ be a diagram in a category $\mathcal{C}$.
Assume the index category $\mathcal{I}$ is filtered. In this case
there are three successively stronger notions which pick out an object
$X$ of $\mathcal{C}$. The first is just
$$
X = \colim_{i \in \mathcal{I}} M_i.
$$
Then $X$ comes equipped with the coprojections $M_i \to X$.
A stronger condition would be to require that $X$ is the colimit and
that there exist an $i \in \mathcal{I}$ and a morphism $X \to M_i$ such
that the composition $X \to M_i \to X$ is $\text{id}_X$. An even
stronger condition is the following.

\begin{definition}
\label{definition-essentially-constant-diagram}
Let $M : \mathcal{I} \to \mathcal{C}$ be a diagram in a category
$\mathcal{C}$.
\begin{enumerate}
\item Assume the index category $\mathcal{I}$ is filtered and
let $(X, \{M_i \to X\}_i)$ be a cocone for $M$, see
Remark \ref{remark-cones-and-cocones}. We say $M$ is
{\it essentially constant} with {\it value} $X$ if there exist an
$i \in \mathcal{I}$ and a morphism $X \to M_i$ such that
\begin{enumerate}
\item $X \to M_i \to X$ is $\text{id}_X$, and
\item for all $j$ there exist $k$ and morphisms $i \to k$ and $j \to k$
such that the morphism $M_j \to M_k$ equals the composition
$M_j \to X \to M_i \to M_k$.
\end{enumerate}
\item Assume the index category $\mathcal{I}$ is cofiltered and let
$(X, \{X \to M_i\}_i)$ be a cone for $M$, see
Remark \ref{remark-cones-and-cocones}. We say
$M$ is {\it essentially constant} with {\it value} $X$ if
there exist an $i \in \mathcal{I}$
and a morphism $M_i \to X$ such that
\begin{enumerate}
\item $X \to M_i \to X$ is $\text{id}_X$, and
\item for all $j$ there exist $k$ and morphisms $k \to i$ and $k \to j$
such that the morphism $M_k \to M_j$ equals the composition
$M_k \to M_i \to X \to M_j$.
\end{enumerate}
\end{enumerate}
Please keep in mind Lemma \ref{lemma-essentially-constant-is-limit-colimit}
when using this definition.
\end{definition}

\noindent
Which of the two versions is meant will be clear from context. If there is
any confusion we will distinguish between these by saying that the first
version means $M$ is essentially constant as an {\it ind-object}, and in
the second case we will say it is essentially constant as a {\it pro-object}.
This terminology is further explained in
Remarks \ref{remark-ind-category} and \ref{remark-pro-category}.
In fact we will often use the terminology ``essentially constant system''
which formally speaking is only defined for systems over directed sets.

\begin{definition}
\label{definition-essentially-constant-system}
Let $\mathcal{C}$ be a category. A directed system
$(M_i, f_{ii'})$ is an {\it essentially constant system}
if $M$ viewed as a functor $I \to \mathcal{C}$
defines an essentially constant diagram. A directed inverse system
$(M_i, f_{ii'})$ is an {\it essentially constant inverse system} if
$M$ viewed as a functor $I^{opp} \to \mathcal{C}$ defines an
essentially constant inverse diagram.
\end{definition}

\noindent
If $(M_i, f_{ii'})$ is an essentially constant system and the morphisms
$f_{ii'}$ are monomorphisms, then for all $i \leq i'$ sufficiently large the
morphisms $f_{ii'}$ are isomorphisms. On the other hand, consider the system
$$
\mathbf{Z}^2 \to \mathbf{Z}^2 \to \mathbf{Z}^2 \to \ldots
$$
with maps given by $(a, b) \mapsto (a + b, 0)$. This system is essentially
constant with value $\mathbf{Z}$ but every transition map has a kernel.

\medskip\noindent
Here is an example of a system which is not essentially constant. Let
$M = \bigoplus_{n \geq 0} \mathbf{Z}$ and to let $S : M \to M$ be the
shift operator $(a_0, a_1, \ldots) \mapsto (a_1, a_2, \ldots)$. In this
case the system $M \to M \to M \to \ldots$ with transition maps $S$
has colimit $0$ and the composition $0 \to M \to 0$ is the identity,
but the system is not essentially constant.

\medskip\noindent
The following lemma is a sanity check.

\begin{lemma}
\label{lemma-essentially-constant-is-limit-colimit}
Let $M : \mathcal{I} \to \mathcal{C}$ be a diagram.
If $\mathcal{I}$ is filtered and $M$ is essentially
constant as an ind-object, then $X = \colim M_i$ exists and $M$
is essentially constant with value $X$.
If $\mathcal{I}$ is cofiltered and $M$ is essentially
constant as a pro-object, then $X = \lim M_i$ exists and $M$ is
essentially constant with value $X$.
\end{lemma}

\begin{proof}
Omitted. This is a good exercise in the definitions.
\end{proof}

\begin{remark}
\label{remark-ind-category}
Let $\mathcal{C}$ be a category. There exists a big category
$\text{Ind-}\mathcal{C}$ of {\it ind-objects of} $\mathcal{C}$.
Namely, if $F : \mathcal{I} \to \mathcal{C}$ and
$G : \mathcal{J} \to \mathcal{C}$ are filtered diagrams in $\mathcal{C}$,
then we can define
$$
\Mor_{\text{Ind-}\mathcal{C}}(F, G) =
\lim_i \colim_j \Mor_\mathcal{C}(F(i), G(j)).
$$
There is a canonical functor $\mathcal{C} \to \text{Ind-}\mathcal{C}$
which maps $X$ to the {\it constant system} on $X$. This is a fully
faithful embedding. In this language one sees that a diagram $F$ is
essentially constant if and only if $F$ is isomorphic to a constant system.
If we ever need this material, then we will formulate this into a lemma
and prove it here.
\end{remark}

\begin{remark}
\label{remark-pro-category}
Let $\mathcal{C}$ be a category. There exists a big category
$\text{Pro-}\mathcal{C}$ of {\it pro-objects} of $\mathcal{C}$.
Namely, if $F : \mathcal{I} \to \mathcal{C}$ and
$G : \mathcal{J} \to \mathcal{C}$ are cofiltered diagrams in $\mathcal{C}$,
then we can define
$$
\Mor_{\text{Pro-}\mathcal{C}}(F, G) =
\lim_j \colim_i \Mor_\mathcal{C}(F(i), G(j)).
$$
There is a canonical functor $\mathcal{C} \to \text{Pro-}\mathcal{C}$
which maps $X$ to the {\it constant system} on $X$. This is a fully
faithful embedding. In this language one sees that a diagram $F$ is
essentially constant if and only if $F$ is isomorphic to a constant system.
If we ever need this material, then we will formulate this into a lemma
and prove it here.
\end{remark}

\begin{example}
\label{example-pro-morphism-inverse-systems}
Let $\mathcal{C}$ be a category. Let $(X_n)$ and $(Y_n)$ be inverse
systems in $\mathcal{C}$ over $\mathbf{N}$ with the usual ordering.
Picture:
$$
\ldots \to X_3 \to X_2 \to X_1
\quad\text{and}\quad
\ldots \to Y_3 \to Y_2 \to Y_1
$$
Let $a : (X_n) \to (Y_n)$ be a morphism of pro-objects of $\mathcal{C}$.
What does $a$ amount to? Well, for each $n \in \mathbf{N}$ there should
exist an $m(n)$ and a morphism $a_n : X_{m(n)} \to Y_n$. These morphisms
ought to agree in the following sense: for all $n' \geq n$ there exists an
$m(n', n) \geq m(n'), m(n)$ such that the diagram
$$
\xymatrix{
X_{m(n, n')} \ar[rr] \ar[d] & & X_{m(n)} \ar[d]^{a_n} \\
X_{m(n')} \ar[r]^{a_{n'}} & Y_{n'} \ar[r] & Y_n
}
$$
commutes. After replacing $m(n)$ by $\max_{k, l \leq n}\{m(n, k), m(k, l)\}$
we see that we obtain $\ldots \geq m(3) \geq m(2) \geq m(1)$ and a commutative
diagram
$$
\xymatrix{
\ldots \ar[r] &
X_{m(3)} \ar[d]^{a_3} \ar[r] &
X_{m(2)} \ar[d]^{a_2} \ar[r] &
X_{m(1)} \ar[d]^{a_1} \\
\ldots \ar[r] &
Y_3 \ar[r] &
Y_2 \ar[r] &
Y_1
}
$$
Given an increasing map $m' : \mathbf{N} \to \mathbf{N}$ with $m' \geq m$
and setting $a'_i : X_{m'(i)} \to X_{m(i)} \to Y_i$ the pair
$(m', a')$ defines the same morphism of pro-systems. Conversely, given
two pairs $(m_1, a_1)$ and $(m_1, a_2)$ as above then these define the same
morphism of pro-objects if and only if we can find $m' \geq m_1, m_2$
such that $a'_1 = a'_2$.
\end{example}

\begin{remark}
\label{remark-pro-category-copresheaves}
Let $\mathcal{C}$ be a category. Let $F : \mathcal{I} \to \mathcal{C}$ and
$G : \mathcal{J} \to \mathcal{C}$ be cofiltered diagrams in $\mathcal{C}$.
Consider the functors $A, B : \mathcal{C} \to \textit{Sets}$ defined by
$$
A(X) = \colim_i \Mor_\mathcal{C}(F(i), X)
\quad\text{and}\quad
B(X) = \colim_j \Mor_\mathcal{C}(G(j), X)
$$
We claim that a morphism of pro-systems from $F$ to $G$ is the same thing
as a transformation of functors $t : B \to A$. Namely, given $t$
we can apply $t$ to the class of $\text{id}_{G(j)}$ in $B(G(j))$
to get a compatible system of elements
$\xi_j \in A(G(j)) = \colim_i \Mor_\mathcal{C}(F(i), G(j))$
which is exactly our definition of a morphism in $\text{Pro-}\mathcal{C}$ in
Remark \ref{remark-pro-category}. We omit the construction of a
transformation $B \to A$ given a morphism of pro-objects from $F$ to $G$.
\end{remark}

\begin{lemma}
\label{lemma-image-essentially-constant}
Let $\mathcal{C}$ be a category. Let $M : \mathcal{I} \to \mathcal{C}$
be a diagram with filtered (resp.\ cofiltered) index category $\mathcal{I}$.
Let $F : \mathcal{C} \to \mathcal{D}$ be a functor.
If $M$ is essentially constant as an ind-object (resp.\ pro-object),
then so is $F \circ M : \mathcal{I} \to \mathcal{D}$.
\end{lemma}

\begin{proof}
If $X$ is a value for $M$, then it follows immediately from the
definition that $F(X)$ is a value for $F \circ M$.
\end{proof}

\begin{lemma}
\label{lemma-characterize-essentially-constant-ind}
Let $\mathcal{C}$ be a category. Let $M : \mathcal{I} \to \mathcal{C}$
be a diagram with filtered index category $\mathcal{I}$.
The following are equivalent
\begin{enumerate}
\item $M$ is an essentially constant ind-object,
\item there exists a cocone $(X, \{M_i \to X\}_i)$ such that for any
$W$ in $\mathcal{C}$ the map
$\colim_i \Mor_\mathcal{C}(W, M_i) \to \Mor_\mathcal{C}(W, X)$
is bijective,
\item $X = \colim_i M_i$ exists and for any $W$ in $\mathcal{C}$
the map
$\colim_i \Mor_\mathcal{C}(W, M_i) \to \Mor_\mathcal{C}(W, X)$
is bijective, and
\item there exists an $i$ in $\mathcal{I}$ and a morphism $X \to M_i$
such that for any $W$ in $\mathcal{C}$ the map $\Mor_\mathcal{C}(W, X) \to
\colim_{j \in \mathcal{I}} \Mor_\mathcal{C}(W, M_j)$
is bijective.
\end{enumerate}
In cases (2), (3), and (4) the value of the essentially constant system is $X$.
\end{lemma}

\begin{proof}
It is clear that (3) implies (2). Assume (2).
Then $\text{id}_X \in \Mor_\mathcal{C}(X, X)$
comes from a morphism $X \to M_i$ for some $i$, i.e., $X \to M_i \to X$
is the identity. Then both maps
$$
\Mor_\mathcal{C}(W, X)
\longrightarrow
\colim_i \Mor_\mathcal{C}(W, M_i)
\longrightarrow
\Mor_\mathcal{C}(W, X)
$$
are bijective for all $W$ where the first one is induced by the morphism
$X \to M_i$ we found above, and the composition is the identity. This means
that the composition
$$
\colim_i \Mor_\mathcal{C}(W, M_i)
\longrightarrow
\Mor_\mathcal{C}(W, X)
\longrightarrow
\colim_i \Mor_\mathcal{C}(W, M_i)
$$
is the identity too. Setting $W = M_j$ and starting with $\text{id}_{M_j}$
in the colimit, we see that $M_j \to X \to M_i \to M_k$ is equal to
$M_j \to M_k$ for some $k$ large enough. This proves (1) holds.

\medskip\noindent
Assume (4). Let $k$ be an object of $\mathcal{I}$. Setting $W = M_k$
we deduce there exists a unique morphism $M_k \to X$ such that
there exists a $j$ and morphisms $k \to j$ and
$i \to j$ in $\mathcal{I}$ such that $M_k \to X \to M_i \to M_j$
is equal to $M_k \to M_j$. The uniqueness guarantees that we obtain a
cocone $(X, \{M_k \to X\})$. In this way we see that (4) implies (2);
some details omitted.

\medskip\noindent
We omit the proof that (1) implies the other conditions.
\end{proof}

\begin{lemma}
\label{lemma-characterize-essentially-constant-pro}
Let $\mathcal{C}$ be a category. Let $M : \mathcal{I} \to \mathcal{C}$
be a diagram with cofiltered index category $\mathcal{I}$.
The following are equivalent
\begin{enumerate}
\item $M$ is an essentially constant pro-object,
\item there exists a cone $(X, \{X \to M_i\})$
such that for any $W$ in $\mathcal{C}$ the map
$\colim_{i \in \mathcal{I}^{opp}} \Mor_\mathcal{C}(M_i, W) \to
\Mor_\mathcal{C}(X, W)$ is bijective,
\item $X = \lim_i M_i$ exists and for any $W$ in $\mathcal{C}$
the map
$\colim_{i \in \mathcal{I}^{opp}} \Mor_\mathcal{C}(M_i, W) \to
\Mor_\mathcal{C}(X, W)$ is bijective, and
\item there exists an $i$ in $\mathcal{I}$ and a morphism $M_i \to X$
such that for any $W$ in $\mathcal{C}$ the map $\Mor_\mathcal{C}(X, W) \to
\colim_{j \in \mathcal{I}^{opp}} \Mor_\mathcal{C}(M_j, W)$ is bijective.
\end{enumerate}
In cases (2), (3), and (4) the value of the essentially constant system is $X$.
\end{lemma}

\begin{proof}
This lemma is dual to Lemma \ref{lemma-characterize-essentially-constant-ind}.
\end{proof}

\begin{lemma}
\label{lemma-cofinal-essentially-constant}
Let $\mathcal{C}$ be a category. Let $H : \mathcal{I} \to \mathcal{J}$
be a functor of filtered index categories. If $H$ is cofinal, then
any diagram $M : \mathcal{J} \to \mathcal{C}$ is essentially constant
if and only if $M \circ H$ is essentially constant.
\end{lemma}

\begin{proof}
This follows formally from
Lemmas \ref{lemma-characterize-essentially-constant-ind} and
\ref{lemma-cofinal}.
\end{proof}

\begin{lemma}
\label{lemma-essentially-constant-over-product}
Let $\mathcal{I}$ and $\mathcal{J}$ be filtered categories and denote
$p : \mathcal{I} \times \mathcal{J} \to \mathcal{J}$ the projection.
Then $\mathcal{I} \times \mathcal{J}$ is filtered and a diagram
$M : \mathcal{J} \to \mathcal{C}$ is essentially constant if and only
if $M \circ p : \mathcal{I} \times \mathcal{J} \to \mathcal{C}$
is essentially constant.
\end{lemma}

\begin{proof}
We omit the verification that $\mathcal{I} \times \mathcal{J}$ is
filtered. The equivalence follows from
Lemma \ref{lemma-cofinal-essentially-constant}
because $p$ is cofinal (verification omitted).
\end{proof}

\begin{lemma}
\label{lemma-initial-essentially-constant}
Let $\mathcal{C}$ be a category. Let $H : \mathcal{I} \to \mathcal{J}$
be a functor of cofiltered index categories. If $H$ is initial, then
any diagram $M : \mathcal{J} \to \mathcal{C}$ is essentially constant
if and only if $M \circ H$ is essentially constant.
\end{lemma}

\begin{proof}
This follows formally from
Lemmas \ref{lemma-characterize-essentially-constant-pro},
\ref{lemma-initial}, \ref{lemma-cofinal}, and
the fact that if $\mathcal{I}$ is initial in $\mathcal{J}$,
then $\mathcal{I}^{opp}$ is cofinal in $\mathcal{J}^{opp}$.
\end{proof}





